%%%%%%%%%%%%%%%%%%%%%%%%%%%%%%%%%%%%%%%%%%%%%%%%%%%%%%%%%%%%%%%%%%%%%%
% Lớp 11 | Các số đặc trưng đo xu thế trung tâm cho mẫu số liệu | Trung vị và tứ phân vị của mẫu số liệu ghép nhóm
% Dạng 1D5?2-1: Câu hỏi lý thuyết
% Số câu: 41  (sinh tự động bởi tach_cau.py; từ nay thầy sửa trực tiếp file này)
%%%%%%%%%%%%%%%%%%%%%%%%%%%%%%%%%%%%%%%%%%%%%%%%%%%%%%%%%%%%%%%%%%%%%%

%%%%%%%%%%%%%%%% TRẮC NGHIỆM 4 PHƯƠNG ÁN %%%%%%%%%%%%%%%%

% nguồn: Ok_TracNghiemLop11_Lan1
\begin{ex}%[1D5N2-1]%[ma:f2c202f9]
	Khi độ chênh lệch các số liệu trong mẫu quá lớn thì đại lượng  nào thích hợp đại diện cho các số liệu trong mẫu.
	 \choice
		{Số trung  bình}
		{\True Số trung  vị}
		{Phương sai} 
		 {Độ lệch chuẩn}
	 \loigiai{
		Khi độ chênh lệch các số liệu trong mẫu quá lớn thì Số trung vị là đại lượng thích hợp  đại diện cho các số liệu trong mẫu.
	} 
\end{ex}

% nguồn: Ok_TracNghiemLop11_Lan1
\begin{ex}%[1D5B2-1]%[ma:cc7b6efc]
	Đạo  hàm của hàm số $y=\sqrt{3x^2+4}$ là
	\choice 
		{$y'=\dfrac{1}{2\sqrt{3x^2+4}}$} 
		{$y'=\dfrac{x}{\sqrt{3x^2+4}}$} 
		{$y'=\dfrac{6x}{\sqrt{3x^2+4}}$} 
		 {\True $y'=\dfrac{3x}{\sqrt{3x^2+4}}$}
	\loigiai{ 
		Ta có $y'=\dfrac{(3x^2+4)'}{2\sqrt{3x^2+4}}=\dfrac{6x}{2\sqrt{3x^2+4}}=\dfrac{3x}{\sqrt{3x^2+4}}$. 
	} 
\end{ex}

% nguồn: Ok_TracNghiemLop11_Lan1
\begin{ex}%[1D5B2-1]%[ma:9ad054b2]
	Cho hàm số $y=\sqrt{3x-1}$. Giá trị  của $y^3\cdot y''$ bằng
	\choice 
		{\True  $-\dfrac{9}{4} $}
		 {$\dfrac{3}{2} $}
		{$-1 $} 
		{$-\dfrac{3}{2}  $}
	\loigiai{ 
		Ta có $y'=\dfrac{3}{2\sqrt{3x-1}}$, $y''=-\dfrac{3}{2} \cdot  \dfrac{\left(\sqrt{3x-1}\right)'}{3x-1} =-\dfrac{9}{4}\cdot \dfrac{1}{(3x-1)\sqrt{3x-1}}$.\\
		Khi  đó $y^3 \cdot y'' =-\dfrac{9}{4}$.
	} 
\end{ex}

% nguồn: Ok_TracNghiemLop11_Lan1
\begin{ex}%[1D5B2-1]%[ma:67956888]
	Đạo hàm của hàm số $y=\cos 3 x$  tại $x=\dfrac{\pi}{2}$ là
	 \choice
		{\True  $3$}
		 {$0$}
		{$1$} 
		{$-3$} 
	 \loigiai{
		 Ta có $y'=-3\sin 3x$, nên $y'\left(\dfrac{\pi}{2}\right)=-3\sin \dfrac{3\pi}{2}=3$.
	 }
\end{ex}

% nguồn: Ok_TracNghiemLop11_Lan1
\begin{ex}%[1D5K2-1]%[ma:eeeaed18]
	Biết đạo hàm của hàm số $y=\dfrac{x-1}{\sqrt{x^{2}+1}}$ là $y'=\dfrac{ax+b}{\sqrt{\left(x^{2}+1\right)^{c}}}$ với $a$, $b$, $c$ là các số nguyên dương. Khi đó giá trị của $2a+b+c$ bằng: 
	\choice 
		 {$5$}
		 {\True $6$}
		{$7$} 
		{$4$} 
	 \loigiai{
		Ta có $y\cdot \sqrt{x^2+1} = x-1$.\\ 
		Lấy đạo hàm hai vế ta có $y'\sqrt{x^2+1}+y\cdot \left(\sqrt{x^2+1}\right)'=1 \Leftrightarrow y'\sqrt{x^2+1}+y\cdot  \dfrac{x}{\sqrt{x^2+1}}=1$.\\
		Suy ra  $y'\sqrt{x^2+1}=1-y\cdot \dfrac{x}{\sqrt{x^2+1}}=1-\dfrac{x^2-x}{x^2+1} = \dfrac{x+1}{x^2+1}$\\
		$\Rightarrow y'  = \dfrac{x}{(x^2+1)\sqrt{x^2+1}}=\dfrac{x}{\sqrt{(x^2+1)^3}}$.\\
		Vậy $a=1,b=0,c=3  \Rightarrow 2a+b+c=6$.
	 }
\end{ex}

% nguồn: Ok_TracNghiemLop11_Lan1
\begin{ex}%[1D5B2-1]%[ma:a16c5396]
	Cho hàm số $f(x)=(2 x-1)^3$. Giá trị của $f"(1)$  bằng
	\choice 
		{$12$} 
		{$6$} 
		{\True  $24$}
		 {$4$}
	 \loigiai{
		 Ta có \\
		$\bullet\,\, f'(x)=3\cdot  (2 x-1)^2\cdot (2x-1)'=6 (2 x-1)^2$.\\
		 $\bullet\,\, f"(x)=12 (2x-1)(2x-1)'=24 (2x-1).$\\
		Vậy $f"(1)=24  (2\cdot 1-1)=24.$
	} 
\end{ex}

% nguồn: Ok_TracNghiemLop11_Lan1
\begin{ex}%[1D5Y2-1]%[ma:bb6e8bfb]
	Cho $u=u(x), v=v(x)$ là các hàm số có  đạo hàm tại điểm $x$ thuộc khoảng xác định. Mệnh đề nào sau đây đúng?
	\choice 
		{\True  $(uv)' = u'v+uv'$}
		{$(uv)'=u'v'$ } 
		 {$\left(\dfrac{u}{v}\right)' = \dfrac{u'}{v'}$ }
		{$\left(\dfrac{u}{v}\right)'  = \dfrac{u'v+uv'}{v^2}$ }
	 \loigiai{
		 Ta có $(uv)'=u'v+uv'; \left(\dfrac{u}{v}\right)' = \dfrac{u'v-uv'}{v^2}$.
	 }
\end{ex}

% nguồn: Ok_TracNghiemLop11_Lan1
\begin{ex}%[1D5Y2-1]%[ma:fbcf775a]
	Cho hàm số $f(x)=\left( 2x-3 \right)^{\frac{5}{6}}$. Tính  $f'(2)$
	 \choice
		{$  \dfrac{5}{6} $}
		{\True  $ \dfrac{5}{3} $}
		{$  -\dfrac{5}{6} $}
		{$ -\dfrac{5}{3}  $}
	\loigiai{ 
		Ta có $f'(x)=\dfrac{5}{6} \cdot 2 \cdot \left( 2x-3 \right)^{-\frac{1}{6}}$ suy ra $f'(2)=  \dfrac{5}{3}$.
	} 
\end{ex}

% nguồn: Ok_TracNghiemLop11_Lan1
\begin{ex}%[1D5B2-1]%[ma:2363cb5b]
	Cho hàm số $y=x^2-2x+2011$. Phương trình $y'=0$  có nghiệm là
	\choice 
		{$x=2$} 
		{\True  $x=1$}
		{$x=-1$} 
		 {$x=-2$}
	\loigiai{ 
		Ta có $y'=0\Leftrightarrow 2x-2=0\Leftrightarrow x=1$. 
	 }
\end{ex}

% nguồn: Ok_TracNghiemLop11_Lan1
\begin{ex}%[1D5B2-1]%[ma:b25e0db2]
	Hàm số $y = \sin \left( {\dfrac{\pi }{2}  - 2x} \right)$ có đạo hàm bằng \vspace{3pt}
	\choice 
		{$y' = 2\sin  2x$}
		{\True $y'  = - 2\cos 2x$}
		{$y' = \cos \left( {\dfrac{\pi  }{2} - 2x} \right)$}
		{$y' = 2\cos \left( {\dfrac{\pi }{2} - 2x} \right)$} 
	\loigiai{Ta  có $$y' =\left( {\dfrac{\pi }{2} - 2x} \right)'\cdot \cos \left( {\dfrac{\pi }{2} - 2x} \right)=- 2\cos 2x .$$
	} 
\end{ex}

% nguồn: Ok_TracNghiemLop11_Lan1
\begin{ex}%[1D5K2-1]%[ma:8e27aa00]
	Cho hàm số $y=f(x)$ có đạo hàm trên $\mathbb{R}$. Xét các hàm số $g(x)=f(x)-f(2x)$ và  $h(x)=f(x)-f(4x)$. Biết rằng $g{'}(1)=18$ và $g{'}(2)=1000$. Tính $h{'}(1)$.
	 \choice
		{$-2018$} 
		{\True  $2018$}
		{$2020$} 
		{$-2020$} 
	\loigiai{ 
		Từ giả thiết  $g(x)=f(x)-f(2x)\Rightarrow g{'}(x)=f{'}(x)-2f{'}(2x)$.\\
		Mặt khác $\heva{g{'}(1)&=18\\g{'}(2)&=1000}\Leftrightarrow \heva{&f{'}(1)-2f{'}(2)=18\\&f{'}(2)-2f{'}(4)=1000}\Leftrightarrow  \heva{&f{'}(1)-2f{'}(2)=18\\&2f{'}(2)-4f{'}(4)=2000}\\\Rightarrow f{'}(1)-4f{'}(4)=2018$.\\
		Lại  có $h(x)=f(x)-f(4x)\Rightarrow h{'}(x)=f{'}(x)-4f{'}(4x)\Rightarrow h{'}(1)=f{'}(1)-4f{'}(4)=2018.$
	} 
\end{ex}

% nguồn: Ok_TracNghiemLop11_Lan1
\begin{ex}%[1D5B2-1]%[ma:0bcd9141]
	Đạo hàm của hàm số $f(x)=5x^3-x^2-1$ trên khoảng $(-\infty;+\infty)$ là 
	 \choice
		{$0$} 
		{$15x^2-2x-1$} 
		 {$15x^2+2x$}
		{\True $15x^2-2x$} 
	\loigiai{Ta  có
		$f'(x)=5\cdot  3x^2-2x=15x^2-2x$.
	} 
\end{ex}

% nguồn: Ok_TracNghiemLop11_Lan1
\begin{ex}%[1D5B2-1]%[ma:24048ead]
	Cho hàm số $f(x)=x^3-2x^2+x-4$. Biết tập nghiệm của bất phương  trình $f'(x)\le 0$ là đoạn $[a;b]$. Tính $P=3a-4b$.
	\choice 
		 {\True $ -3 $}
		{$ -1  $}
		{$  25 $}
		{$  -\dfrac{5}{3} $}
	\loigiai{ 
		Ta có $f'(x)=3x^2-4x+1$.\\ 
		$f'(x)\le  0 \Leftrightarrow x\in \left[\dfrac{1}{3};1\right]$.\\
		Khi đó $a=\dfrac{1}{3}$,  $b=1$.\\
		 Vậy $P=3a-4b=1-4=-3$.
	} 
\end{ex}

% nguồn: Ok_TracNghiemLop11_Lan1
\begin{ex}%[1D5K2-1]%[ma:cd1a7c0c]
	Cho hàm số $y=(m+2)x^3+\dfrac{3}{2}(m+2)x^2+3x-1$, $m$ là tham số. Số các giá trị nguyên $m$ để $y'\ge 0$, $\forall x\in \mathbb{R}$  là
	\choice 
		{\True  $5$}
		{$3$} 
		{Có vô số giá trị nguyên  $m$}
		{$4$} 
	\loigiai{ 
		Ta có  $y'=3(m+2)x^2+3(m+2)x+3$.\\
		{\bf TH1:} $m+2=0\Leftrightarrow m=-2 \Rightarrow y'=3\ge 0$, $\forall x\in \mathbb{R}$.\\ 
		{\bf TH2:} $m+2\ne0\Leftrightarrow m\ne-2$. 
		$$y'\ge 0,\ \forall x\in \mathbb{R} \Leftrightarrow \heva{& m+2>0 \\ & \Delta  \le 0}\Leftrightarrow \heva{& m>-2 \\ & (m+2)(m-2) \le 0}\Leftrightarrow \heva{& m>-2 \\ &m\in [-2;2]}\Leftrightarrow m\in (-2;2].$$
		Vì $m\in \mathbb{Z}$ nên $m\in  \{-2,-1,0,1,2\}$.Vậy có 5 giá trị nguyên của $m$.
	} 
\end{ex}

% nguồn: Ok_TracNghiemLop11_Lan1
\begin{ex}%[1D5Y2-1]%[ma:0ffd5ed8]
	Cho $ u=u(x), v=v(x), v(x)\ne 0 $.  Hãy chọn khẳng định \textbf{sai}?
	\choice 
		 {$(u+v)'=u'+v'$}
		 {\True $\left(\dfrac{1}{v}\right)'=-\dfrac{v'}{v}$}
		 {$(uv)'=u'v+uv'$}
		{$(ku)'=ku'$} 
	 \loigiai{
		 Ta có $\left(\dfrac{1}{v}\right)'=-\dfrac{v'}{v^2}$.
	} 
\end{ex}

% nguồn: Ok_TracNghiemLop11_Lan1
\begin{ex}%[1D5Y2-1]%[ma:c303e24e]
	Công thức nào  sau đây \textbf{sai}?
	 \choice
		{\True $(C)'=C$ với  $C$ là hằng số}
		 {$\left(x^n\right)'=n\cdot x^{n-1}; n\in\mathbb{N}^*, n>1$}
		{$\left(\sqrt{x}\right)'=\dfrac{1}{2\sqrt{x}};  x>0$}
		{$(x)'=1$} 
	\loigiai{ 
		Ta  có $(C)'=0$.
	 }
\end{ex}

% nguồn: Ok_TracNghiemLop11_Lan1
\begin{ex}%[1D5Y2-1]%[ma:5cae25c3]
	Hàm số nào  dưới đây có đạo hàm $y'=3x^2-6x$?
	 \choice
		{$y=x^3-3x^2+x$} 
		 {$y=x^3+3x^2-2$}
		 {$y=x^3-3x^2+2x$}
		 {\True $y=x^3-3x^2+2$}
	 \loigiai{
		\begin{center}
			$y=x^3-3x^2+2\Rightarrow  y'=3x^2-6x$.
		\end{center}
	} 
\end{ex}

% nguồn: Ok_TracNghiemLop11_Lan1
\begin{ex}%[1D5Y2-1]%[ma:fdb71914]
	Tìm đạo  hàm của hàm số $y=x^4-2x^2-3$.
	 \choice
		 {\True $y'=4x^3-4x$}
		{$y'=x^3-2x$} 
		{$y'=4x^3-4x-3$} 
		{$y'=\dfrac{1}{4}x^3-x$} 
	\loigiai{ 
		 $y'=4x^3-4x$.
	} 
\end{ex}

% nguồn: Ok_TracNghiemLop11_Lan1
\begin{ex}%[1D5B2-1]%[ma:808b8c3e]
	Đạo hàm cấp hai  của hàm số $y=x^3+2x^2+\dfrac{1}{7}$ bằng
	 \choice
		 {$y''=3x$}
		{$y''=3x^2+2x$} 
		{$y''=3x+1$} 
		{\True $y''=6x+4$} 
	\loigiai{ 
		Ta có $y'=3x^2+4x$, $y''=6x+4$. 
	} 
\end{ex}

% nguồn: Ok_TracNghiemLop11_Lan1
\begin{ex}%[1D5Y2-1]%[ma:3ce1b776]
	Đạo hàm của hàm số $y=(1-x)^{2019}$ tại  $x=0$ bằng
	\choice 
		{$2019$} 
		{\True $-2019$} 
		{$-2019\cdot  x^{2018}$}
		{$2019x^{2018}$} 
	 \loigiai{
		Có $y'=-2019\cdot (1-x)^{2018}  \Rightarrow y'(0)=-2019$.
	 }
\end{ex}

% nguồn: Ok_TracNghiemLop11_Lan1
\begin{ex}%[1D5Y2-1]%[ma:aadb2a9f]
	Tìm đạo hàm của hàm số $y =  x^2 +x +1$ có đạo hàm trên $\mathbb{R}$.
	 \choice
		{$y'=2+x$} 
		{\True  $y'=2x+1$}
		{$y'=3x$} 
		{$y'=x^2+x$} 
	 \loigiai{
		Ta có $y'=(x^2+x+1)'=2x+1$. 
	} 
\end{ex}

% nguồn: Ok_TracNghiemLop11_Lan1
\begin{ex}%[1D5Y2-1]%[ma:eef06a6d]
	Đạo hàm của hàm số $y=\left( x^3-2x^2  \right)^2$ bằng
	\choice 
		 {$6x^5-20x^4-16x^3$}
		{$6x^5-20x^4+4x^3$} 
		{$6x^5+16x^3$} 
		{\True  $6x^5-20x^4+16x^3$}
	 \loigiai{
		 Ta có $y'=2\left( x^3-2x^2 \right)(3x^2-4x)=6x^5-20x^4+16x^3$.
	} 
\end{ex}

% nguồn: Ok_TracNghiemLop11_Lan1
\begin{ex}%[1D5G2-1]%[ma:6acb76e3]
	Cho hàm  số $f(x)=\dfrac{4x+50}{x^2-25}$. Tính $f^{(2022)}(4)$ ta được kết quả là
	\choice 
		 {$-7\cdot 2022!-\dfrac{4\cdot 2022!}{9^{2023}}$}
		{$\left(-\dfrac{22}{3}\right)^{2022}$} 
		{$-7\cdot 2022!+\dfrac{3\cdot 2022!}{9^{2023}}$} 
		{\True $-7\cdot  2022!-\dfrac{3\cdot 2022!}{9^{2023}}$}
	 \loigiai{
		Ta  có $f(x)=\dfrac{4x+50}{x^2-25}=\dfrac{7}{x-5}-\dfrac{3}{x+5}$.\\
		Xét hàm số $y=\dfrac{1}{ax+b}$  ta có $y^{(n)}=\dfrac{(-1)^n\cdot a^n\cdot n!}{(ax+b)^{n+1}}$.\\
		Áp dụng  công thức trên ta có $f^{n}(x)=7\cdot \dfrac{(-1)^n\cdot 1^n\cdot n!}{(x-5)^{n+1}}-3\cdot\dfrac{(-1)^n\cdot 1^n\cdot n!}{(x+5)^{n+1}}$.\\
		Vậy $f^{(2022)}(4)=-7\cdot  2022!-\dfrac{3\cdot 2022!}{9^{2023}}$.
	} 
\end{ex}

% nguồn: Ok_TracNghiemLop11_Lan1
\begin{ex}%[1D5B2-1]%[ma:d60a42eb]
	Cho hai hàm số  $f(x)=x^2+5$ và $g(x)=9x-\dfrac{3}{2}x^2$. Giá trị của $x$ là bao nhiêu để $f'(x)=g'(x)$?
	 \choice
		 {$-4$}
		 {$4$}
		{\True $\dfrac{9}{5}$} 
		{$\dfrac{5}{4}$} 
	\loigiai{ 
	Ta có $\heva{&f'(x)=2x\\&g'(x)=9-3x}$,  $f'(x)=g'(x) \Leftrightarrow 2x=9-3x \Leftrightarrow x=\dfrac{9}{5}$.
	 }
\end{ex}

% nguồn: Ok_TracNghiemLop11_Lan1
\begin{ex}%[1D5Y2-1]%[ma:e80bb502]
	Cho  hàm số $f(x)=2x^{3}+1$. Giá trị $f'(-1)$ bằng:
	 \choice
		{\True  $6$}
		{$3$} 
		{$-2$} 
		 {$-6$}
	\loigiai{ 
		Ta có $f'(x) =  6x^2 \Rightarrow f'(-1) = 6$.
	 }
\end{ex}

% nguồn: Ok_TracNghiemLop11_Lan1
\begin{ex}%[1D5B2-1]%[ma:6214837b]
	Đạo hàm  của hàm số $y=\sqrt{4x^2+3x+1}$ là
	\choice 
		 {$y'=12x+3$}
		{$y'=\dfrac{1}{2\sqrt{4x^2+3x+1}}$} 
		{\True  $y'=\dfrac{8x+3}{2\sqrt{4x^2+3x+1}}$}
		 {$y'=\dfrac{8x+3}{\sqrt{4x^2+3x+1}}$}
	 \loigiai{
	Ta có  $y'=\dfrac{(4x^2+3x+1)'}{2\sqrt{4x^2+3x+1}}=\dfrac{8x+3}{2\sqrt{4x^2+3x+1}}$.}
\end{ex}

% nguồn: Ok_TracNghiemLop11_Lan1
\begin{ex}%[1D5G2-1]%[ma:fe504e89]
	Tính giá trị của 
	$A=2 \cdot \mathrm{C}_{2020}^1+2\cdot 2^2\cdot \mathrm{C}_{2020}^2+3\cdot 2^3\cdot \mathrm{C}_{2020}^3+\cdots+2020\cdot 2^{2020}\cdot  \mathrm{C}_{2020}^{2020}$.
	 \choice
		 {$A=4040\cdot 3^{2020}$}
		{\True  $A=4040\cdot 3^{2019}$}
		 {$A=2020\cdot 3^{2019}$}
		{$A=2020\cdot  3^{2020}$}
	\loigiai{ 
		Lấy đạo hàm hai vế  của $(1+x)^n=\sum\limits_{k=0}^{n}\mathrm{C}_n^k x^k$ ta được
		$$n(1+x)^{n-1}=\sum\limits_{k=0}^{n} k\cdot  \mathrm{C}_n^k x^{k-1}.$$
		Tiếp tục nhân $2$ vế cho $x$, ta được 
		$$nx(1+x)^{n-1}=\sum\limits_{k=0}^{n} k\cdot \mathrm{C}_n^k x^{k}.$$ 
		Chọn $x=2$ và $n=2020$, ta  được $$2020\cdot 2(1+2)^{2019}=\sum\limits_{k=0}^{2020} k\cdot \mathrm{C}_{2020}^k 2^{k}.$$
		Vậy $4040\cdot 3^{2019}=\sum\limits_{k=0}^{2020} k\cdot \mathrm{C}_{2020}^k 2^{k}$ hay $A=4040\cdot  3^{2019}$.
	} 
\end{ex}

% nguồn: Ok_TracNghiemLop11_Lan1
\begin{ex}%[1D5B2-1]%[ma:b211a854]
	Cho  hàm số $f(x)=\sqrt{3+x}$. Tính $f(1)+4f'(1)$.
	\choice 
		{$1$} 
		{\True $3$} 
		 {$\dfrac{1}{4}$}
		 {$0$}
	 \loigiai{
		Ta có $f(x)=\sqrt{3+x} \Rightarrow  f'(x)=\dfrac{1}{2\sqrt{3+x}}$. \\
	Nên  $f(1)+4f'(1)=\sqrt{3+1}+\dfrac{4}{2\sqrt{3+1}}=3$.}
\end{ex}

% nguồn: Ok_TracNghiemLop11_Lan1
\begin{ex}%[1D5Y2-1]%[ma:4e559416]
	Đạo hàm của hàm số $y=\sqrt{x}$  tại điểm $x=9$ bằng
	 \choice
		 {$0$}
		 {$\dfrac{1}{2}$}
		{\True  $\dfrac{1}{6}$}
		{$\dfrac{1}{3}$} 
	\loigiai{ 
		Ta có $y'=\dfrac{1}{2\sqrt{x}}$. Từ  đó $y'(9)=\dfrac{1}{6}$.
	} 
\end{ex}

% nguồn: Ok_TracNghiemLop11_Lan1
\begin{ex}%[1D5Y2-1]%[ma:79e483f6]
	Cho  các hàm số $u=u(x)$, $v=v(x)$ có đạo hàm tại điểm $x$ thuộc khoảng xác định. Khẳng định nào dưới đây {\bf sai}?
	 \choice
		{\True  $\left(\dfrac{u}{v}\right)'=\dfrac{u'v+uv'}{v^2}$, ($v=v(x)\neq 0$)}
		{$(u+v)'=u'+v'$} 
		 {$(uv)'=u'v+uv'$}
		{$(u-v)'=u'-v'$} 
	\loigiai{ 
		Công thức đúng là  $\left(\dfrac{u}{v}\right)'=\dfrac{u'v-uv'}{v^2}$, ($v=v(x)\neq 0$).
	 }
\end{ex}

% nguồn: Ok_TracNghiemLop11_Lan1
\begin{ex}%[1D5K2-1]%[ma:33cd2c5c]
	Hàm số  nào sau đây có đạo hàm luôn dương $\forall x\in \mathbb{R}$?
	\choice 
		{$y=\dfrac{x-1}{x+3}$} 
		{\True $y=2x^3-3x^2+5x+9$} 
		{$y=x^5-x^3$} 
		 {$y=\tan x+x$}
	 \loigiai{
		Xét các  đáp án
		\begin{itemize}
			\item Hàm số $y=\dfrac{x-1}{x+3}$ có tập xác định là $\mathscr{D}=\mathbb{R}\setminus\{-3\}$ và $y'=\dfrac{4}{(x+3)^2}$ nên đạo hàm của hàm số không thể luôn dương $\forall x\in \mathbb{R}$. 
			\item Hàm số $y=2x^3-3x^2+5x+9$ có tập xác định là $\mathscr{D}=\mathbb{R}$ và $y'=6x^2-6x+5=6\left(x-\dfrac{1}{2}\right)^2+\dfrac{7}{2}>0,\ \forall x\in \mathbb{R}$ nên đạo hàm của  hàm số luôn dương $\forall x\in \mathbb{R}$.
			\item Hàm số $y=x^5-x^3$ có tập xác định là $\mathscr{D}=\mathbb{R}$; $y'=5x^4-3x^2$ và $y'=0$ có nghiệm nên đạo hàm của hàm  số không thể luôn dương $\forall x\in \mathbb{R}$.
			\item Hàm số $y=\tan x+x$ có tập xác định là $\mathscr{D}=\mathbb{R}\setminus \left\{\dfrac{\pi}{2}+k\pi,\ k\in \mathbb{Z} \right\}$ và  $y'=\dfrac{1}{\cos^2 x}+1$ nên đạo hàm của hàm số không thể luôn dương $\forall x\in \mathbb{R}$.
		\end{itemize}
	} 
\end{ex}

% nguồn: Ok_TracNghiemLop11_Lan1
\begin{ex}%[1D5Y2-1]%[ma:e527a723]
	Cho  hàm số $ f(x)=x^4-3x^2+5. $ Tính $ f'(2).$
	 \choice
		{$ 4  $}
		{$  5 $}
		{\True $ 20  $}
		{$ 0  $}
	\loigiai{ 
		 Ta có $ f'(x)=4x^3-6x\Rightarrow f'(2)=20. $
	} 
\end{ex}

% nguồn: Ok_TracNghiemLop11_Lan1
\begin{ex}%[1D5Y2-1]%[ma:aa6207f6]
	Cho hàm số $y=\cos 2x$. Công thức nào sau đây là  đúng?
	\choice 
		 {$y'=-\sin 2x$}
		{\True  $y'=-2\sin 2x$}
		{$y'=\sin 2x$} 
		{$y'=2\sin  2x$}
	 \loigiai{
		Ta có  $y'=(\cos 2x)'=-(2x)'\sin 2x=-2\sin 2x$.
	 }
\end{ex}

% nguồn: Ok_TracNghiemLop11_Lan1
\begin{ex}%[1D5Y2-1]%[ma:aabdd9f8]
	Hàm số $y=\sqrt{x^4+1}$ có đạo  hàm $y'$ bằng
	\choice 
		{$\dfrac{1}{\sqrt{x^4+1}}$} 
		{$\dfrac{4x^3}{\sqrt{x^4+1}}$} 
		{\True  $\dfrac{2x^3}{\sqrt{x^4+1}}$}
		 {$\dfrac{x^4}{2\sqrt{x^4+1}}$}
	 \loigiai{
	Ta  có $y'=\dfrac{(x^4+1)'}{2\sqrt{x^4+1}}=\dfrac{4x^3}{2\sqrt{x^4+1}}=\dfrac{2x^3}{\sqrt{x^4+1}}$.
	 }
\end{ex}

% nguồn: Ok_TracNghiemLop11_Lan1
\begin{ex}%[1D5Y2-1]%[ma:7b073fa3]
	Cho các hàm số $u=u(x), v=v(x)$ có đạo hàm tại mọi $x$ thuộc tập xác định. Khẳng định nào sau đây là  \textbf{sai}?
	\choice 
		{$(u+v)'=u'+v'$} 
		{$(uv)'=u'v+uv'$} 
		{\True $(u^n)'=nu^{n-1}$} 
		{$(\sqrt{u})'=\dfrac{u'}{2\sqrt{u}}$} 
	\loigiai{ 
		Khẳng định $(u^n)'=nu^{n-1}$ là sai  vì theo công thức tính đạo hàm của hàm số hợp $(u^n)'=nu^{n-1}u'$.
	 }
\end{ex}

% nguồn: Ok_TracNghiemLop11_Lan1
\begin{ex}%[1D5K2-1]%[ma:b4e4ab11]
	Cho hàm số $y=\dfrac{1}{3}x^3-mx^2+(m+6)x+3$ có đạo hàm là $y'$.  Tìm tất cả các giá trị của $m$ để phương trình $y'=0$ có hai nghiệm phân biệt là $x_1$, $x_2$ thoả mãn $x_1^2+x_2^2=30$.
	 \choice
		{$m=-2$; $m=\dfrac{7}{2}$} 
		{\True  $m=-3$; $m=\dfrac{7}{2}$}
		 {$m=-3$}
		 {$m=3$}
	 \loigiai{
		Ta có  $y'=x^2-2mx+(m+6)$.\\
		Phương trình $y'=0$ có hai nghiệm phân biệt $x_1$,  $x_2$ $\Leftrightarrow \Delta'>0\Leftrightarrow (-m)^2-(m+6)>0\Leftrightarrow \hoac{&m<-2\\&m>3.}$\\
		Theo định lí Viet ta có 
		\[ 
		 \heva{&x_1+x_2=2m\\&x_1\cdot x_2=m+6}
		\] 
		Theo  đề\allowdisplaybreaks
		\begin{eqnarray*}
			 &&x_1^2+x_2^2=30\\
			&\Leftrightarrow&(x_1+x_2)^2-2x_1x_2=30\\ 
			&\Leftrightarrow&(2m)^2-2(m+6)=30\\ 
			&\Leftrightarrow&4m^2-2m-42=0\\ 
			 &\Leftrightarrow&\hoac{&m=\dfrac{7}{2}&(\text{nhận})\\&m=-3&(\text{nhận}).}
		\end{eqnarray*}
	 }
\end{ex}

% nguồn: Ok_TracNghiemLop11_Lan1
\begin{ex}%[1D5B2-1]%[ma:b60469cb]
	Cho hàm số $y=(2x^2+1)^3$. Tìm tập hợp tất cả nghiệm  của bất phương trình $y'\ge 0$.
	\choice 
		{$\mathbb{R}$} 
		{$(-\infty;0]$} 
		 {$\varnothing$}
		{\True $[0;+\infty)$} 
	 \loigiai{
		 Ta có $y'=3(2x^2+1)^2\cdot (2x^2+1)'=12x(2x^2+1)^2$. Ta có
		$$y'\ge 0 \Leftrightarrow 12x(2x^2+1)^2 \ge  0 \Leftrightarrow x\ge 0$$
	} 
\end{ex}

% nguồn: Ok_TracNghiemLop11_Lan1
\begin{ex}%[1D5B2-1]%[ma:7e1cf734]
	Cho hàm số $f(x)=\sin^2{x}-\cos^2{x}-x$. Khi  đó $f'(x)$ bằng:
	\choice 
		 {$1+2\sin{2x}$}
		 {$1-\sin{2x}$}
		 {$-1+\sin{x}\cos{x}$}
		 {\True $-1+2\sin{2x}$}
	\loigiai{ 
		\begin{center}
			 $f(x)=\sin^2{x}-\cos^2{x}-x\Rightarrow f'(x)=2\sin{x}\cdot\cos{x}+2\cos{x}\cdot\sin{x}-1=2\sin{2x}-1$.
		\end{center}
	} 
\end{ex}

% nguồn: Ok_TracNghiemLop11_Lan1
\begin{ex}%[1D5B2-1]%[ma:f7612147]
	Cho hàm số $f(x)$ có $f'(1)=2$. Đạo hàm của hàm số $3 f(x)$  tại điểm $x=1$ bằng
	\choice 
		{$1$} 
		 {$-1$}
		 {\True $6$}
		{$5$} 
	 \loigiai{
		Ta  có $\left.\left[3 f(x)\right]'\right|_{x=1}=3f'(1)=6$.
	 }
\end{ex}

% nguồn: Ok_TracNghiemLop11_Lan1
\begin{ex}%[1D5Y2-1]%[ma:a406ffca]
	Tính đạo hàm  của hàm số $y=x+\cot x-3 $.
	\choice 
		{$y'=1+\dfrac 1{\sin^2x}$} 
		 {\True $y'=1-\dfrac 1{\sin^2x}$}
		{$y'=1-\dfrac 1{\cos^2x}$} 
		{$y'=1+\dfrac  1{\cos^2x}$}
	\loigiai{ 
		Ta  có $y'=(x+\cot x-3)'=1-\dfrac 1{\sin^2x}$.
	 }
\end{ex}


%%%%%%%%%%%%%%%% ĐÚNG/SAI %%%%%%%%%%%%%%%%

% nguồn: Ok_dungsai_lop11_Lan1
\begin{ex}%[1D5H2-1]%[ma:f50b70ab]
 Bạn An khi tra cứu thông tin cho môn học Địa lý, bạn đã tìm được bảng lượng mưa tháng $8$ từ năm $2002$ đến $2022$ tại Trạm quan trắc ở Nha Trang như sau (theo Tổng cục thống kê).  Giáo viên yêu cầu bạn An sau khi tìm được thông tinthì chuyển mẫu số liệu này sang dạng ghép nhóm gồm 4 nhóm và có độ dài bằng nhau.\\[2mm]
 \textit{Bảng số liệu lượng mưa tháng 8 từ năm  2002 đến 2022 tại Trạm quan trắc ở Nha Trang}.
 \begin{tabular}{|c|c|c|c|c|c|c|c|c|c|c|c|}
  \hline 
  Số năm  & $2002$ & $2003$ & $2004$ & $2005$ & $2006$ & $2007$ & $2008$ & $2009$ & $2010$ & $2011$ & $2012$\\
  \hline 
  Lượng mưa (mm) & $38{,}0$  & $25{,}0$ & $49{,}0$ & $11{,}0$ & $68{,}0$ & $51{,}0$ & $80{,}0$ & $40{,}9$ & $62{,}5$ & $132{,}0$ & $29{,}6$\\
   \hline
 \end{tabular}\\[2.5mm]
 \begin{tabular}{|c|c|c|c|c|c|c|c|c|c|c|}
   \hline
  Số năm & $2013$ & $2014$ & $2015$ &  $2016$ & $2017$ & $2018$ & $20019$ & $2020$ & $2021$ & $2022$\\
   \hline
  Lượng mưa (mm) &  $92{,}5$ & $130{,}8$ & $1{,}4$ & $82{,}2$ & $46{,}3$ & $16{,}0$ & $23{,}9$ & $73{,}7$ & $38{,}1$ & $154{,}5$ \\
   \hline
 \end{tabular}\\
  \choiceTF
  {Bạn  An đã chuyển mẫu số liệu trên sang dạng ghép nhóm gồm $4$ nhóm và có độ dài bằng nhau như sau\\
  \begin{tabular}{|c|c|c|c|c|}
   \hline 
   Tổng lượng mưa trong tháng $8$ (mm) &  $[0;40)$ & $[40;80)$ & $[80;120)$ & $[120;160)$\\
   \hline 
   Số năm & $8$ & $8$ & $2$ & $3$  \\
    \hline
 \end{tabular}}
 {Số trung bình và số trung vị của của mẫu số liệu ghép nhóm này lần lượt  là $61{,}9$ và $54{,}5$}
 {Tứ phân vị thứ nhất và tứ phân vị thứ ba của mẫu số liệu ghép nhóm này  là $26{,}6$ và $90$}
 {Mẫu số  liệu ghép nhóm này không có mốt}
  \loigiai{
  \begin{enumerate}
   \item \textbf{Sai.} 
   Bảng số liệu ghép nhóm đúng là\\ 
   \begin{tabular}{|c|c|c|c|c|}
    \hline 
    Tổng lượng mưa  trong tháng $8$ (mm) & $[0;40)$ & $[40;80)$ & $[80;120)$ & $[120;160)$\\
    \hline 
     Số năm & $8$ & $7$ & $3$ & $3$ \\
     \hline
   \end{tabular}
   \item  \textbf{Sai.}\\
   Số trung  bình $\overline{x}=\dfrac{m_1x_1+\ldots+m_4x_4}{n}=\dfrac{8\cdot20+7\cdot60+3\cdot100+3\cdot140}{21}\approx61{,}90$.\\
   Số trung vị nằm  ở nhóm $[40;80)$.
   \begin{eqnarray*}
    M_e&=&a_p+\dfrac{\dfrac{n}{2}-(m_1+\ldots+m_{p-1})}{m_p}\cdot(a_{p+1}-a_p)\\ 
    &=&40+\dfrac{\dfrac{21}{2}-(8)}{7}\cdot(80-40)\approx  54{,}29.
   \end{eqnarray*}
   \item \textbf{Sai.}\\ 
   Tứ phân vị thứ nhất nằm ở  nhóm $[0;40)$.
   \begin{eqnarray*}
     Q_1&=&a_p+\dfrac{\dfrac{n}{4}-(m_1+\ldots+m_{p-1})}{m_p}\cdot(a_{p+1}-a_p)\\
     &=&0+\dfrac{\dfrac{21}{4}-(0)}{8}\cdot(40-0)=26{,}25.
   \end{eqnarray*}
   Tứ phân vị  thứ ba nằm ở nhóm $[80;120)$
   \begin{eqnarray*}
     Q_3&=&a_p+\dfrac{\dfrac{3n}{4}-(m_1+\ldots+m_{p-1})}{m_p}\cdot(a_{p+1}-a_p)\\
    &=&80+\dfrac{\dfrac{3\cdot21}{4}-(8+7)}{3}\cdot(120-80)=90. 
   \end{eqnarray*}
   \item  \textbf{Sai.}
    Mốt của mẫu số liệu ghép nhóm không tồn tại khi tần số của các nhóm số liệu bằng nhau. \\
   Do nhóm $[0;40)$ có tần số lớn nhất nên mốt nằm ở  nhóm $[0;40)$.\\
   Mốt 
   \begin{eqnarray*}
     M_\circ&=&a_j+\dfrac{(m_j-m_{j-1})}{(m_j-m_{j-1})+(m_j-m_{j+1})}\cdot h\\
    &=&0+\dfrac{8-0}{(8-0)+(8-7)}\cdot 40  \approx 35{,}56.
   \end{eqnarray*}
  \end{enumerate}
 } 
\end{ex}
